\documentclass[11pt]{article}
\usepackage[margin=1in]{geometry}
\usepackage{graphicx,booktabs,longtable}
\usepackage[hidelinks]{hyperref}
\title{PRISM: A Statistical Investigation of Facial Phenotype Variation Using FaceKit}
\author{Blinded for review}
\date{}
\begin{document}
\maketitle
\begin{abstract}
We investigate rare-syndrome facial phenotype associations using 120 geometric measurements supplied by FaceKit. The study includes 1,845 age-eligible recorded GMDB patients and an expanded reference of 3,409 FairFace image paths. Forty-eight syndrome cohorts meet the primary minimum size, yielding 5,760 two-sided syndrome--feature tests. A single global Benjamini--Hochberg correction identifies 955 associations at $q<0.05$, including 482 with $|g|\geq0.8$. Recorded-demographic, within-GMDB, adjusted-model, reference-composition and photograph-selection analyses assess the scope and robustness of these associations. Synthetic PDIDB records are maintained separately and contribute no observations to real-patient inference. PRISM is a statistical study built on FaceKit measurements; the findings describe exploratory cohort associations rather than diagnostic performance.
\end{abstract}
\section{Introduction}
Interpretable facial measurements can describe variation across rare-disease image cohorts while remaining sensitive to image acquisition, landmark detection and geometric normalization. FaceKit supplies the feature-extraction methods used in this study \cite{facekit2026}. PRISM investigates statistical associations in those measurements using GMDB \cite{gestaltmatcher} and FairFace \cite{fairface}. The contribution is the cohort design, statistical investigation and evaluation of robustness across comparison specifications. Feature extraction is attributed to FaceKit.
\section{Materials and Methods}
\subsection{Sources, authorization and analytical roles}
Restored real GMDB images were processed locally under mentor research authorization conveyed by the author, subject to PI confirmation of the exact IRB scope. Original GMDB photographs, patient identifiers and per-person measurements are excluded from manuscript materials and public outputs. FairFace is a reference dataset with recorded demographic labels; clinical health and unique-person identities are unverified. The official margin-0.25 archive and both label files passed ZIP integrity and label-to-image correspondence checks. All 11,764 distinct age-category 3--9 candidate paths were screened before sampling.
The primary reference selects up to 500 accepted paths from each of the seven source race-label categories, without replacement: 3,409 paths, comprising 409 Middle Eastern and 500 in each other category. Alternatives use all 5,864 accepted paths or 151 in each of 14 category-by-recorded-sex cells, totaling 2,114. These overlapping specifications assess reference composition and size; they are not independent replications. FairFace annotations are not measurements of biological ancestry or exact age.
\subsection{FaceKit measurements and cohort selection}
All sources use FaceKit revision \texttt{5764bdd4af9360b55a046a16579fb5a890a8182f} and the same checksummed MediaPipe Face Landmarker float16 version-1 model. RGB decoding introduces no additional cropping, resizing, restoration or EXIF rotation. An up-to-two-face pass screens ambiguity, followed by the standard one-face detector for 478 floating-point pixel landmarks and the transformation matrix. FaceKit supplies matrix-based canonicalization, fixed canonical depth and feature-specific normalization. Acceptance requires full pose correction, exactly 120 finite morphology outputs and absolute yaw/pitch/roll at most $15^\circ/15^\circ/10^\circ$. Angular outputs are degrees; other outputs are unitless geometric indices. Supplementary Table S6 records the exact ordered features, formulas, source pin and hashes. The executed software pin is reported directly; an exact paper-execution revision is not independently established.
Recorded GMDB age is year plus month/12, with nonnegative integer components, month 0--11 and $3\leq\mathrm{age}<10$. Missing, invalid and zero-zero ages are excluded. One eligible image per recorded patient is selected by minimum $\max(|yaw|/15,|pitch|/15,|roll|/10)$, then lexical image key. This yields 1,845 recorded patients across 379 labels, including 872 patients in the 48 primary cohorts with at least ten patients. Recorded IDs implement the sampling unit but do not establish undocumented biological independence. FairFace ages remain source categories and distinct paths do not establish distinct people.
\subsection{Primary statistical family}
The primary family comprises 48 syndrome cohorts times 120 measurements, or 5,760 two-sided Mann--Whitney U tests. Feature eligibility requires at least ten finite observations per side and no more than 30\% missingness. All primary cells are estimable. Hedges' $g$ is syndrome minus reference, with correction $J=1-3/[4(n_1+n_2)-9]$. A single global BH correction spans the complete primary family. Marginal percentile intervals use 500 independent group-row bootstrap resamples, seed 2026; ranking stability uses 200 draws. Confidence intervals are not multiplicity-adjusted. Robust analyses trim scaled-MAD outliers independently within groups. Rank AUC is descriptive within-sample separation and does not measure held-out diagnostic performance.
\subsection{Recorded-demographic analyses}
Broad and detailed source-label mappings define separate analysis layers. Analysis A uses reference-only Kruskal--Wallis tests with a 120-feature BH family per layer; reference-group pairwise comparisons have separate within-pair correction. Analysis B compares syndrome and reference observations within recorded groups, with 1,000 effect bootstraps and a global BH family per layer. Analysis C estimates HC3 disease-by-group interactions with a global BH family per layer. The computational support minimum is two disease and twenty reference observations in required cells; estimates from sparse cells remain exploratory. Group matching and HC3 adjustment do not remove dataset/source confounding or establish causal demographic effects.
\subsection{Within-GMDB comparisons, adjustment and robustness}
Each eligible syndrome is compared with all other age-eligible GMDB patients and, separately, with patients from the other eligible syndromes. Each specification has a separate 5,760-test BH family. The adjusted screen retains 49 prespecified syndrome targets and 120 features (5,880 planned models); targeted specificity comparisons retain 186 prespecified anchor pairs (22,320 hypotheses), and two age strata retain 11,760 planned hypotheses. Unsupported cells contribute correction $p=1$, with missing effect estimates and displayed q-values, so support failures do not reduce the planned denominator. These target families are declared separately from the 48-cohort primary family.
Four prespecified age/sex-adjusted within-GMDB models and three reference-covariate models retain separate four- and three-model correction scopes. Reference covariates are recorded sex, broad group and yaw/pitch/roll. Adjusted coefficients are mean differences in normalized feature units rather than Hedges' $g$. Reference alternatives, deterministic alternate-photo selection and 100 random-photo replicates assess specification sensitivity. Reference-feature Spearman clustering is descriptive and does not reduce BH denominators.
Concordance summaries report descriptive point estimates only. The assessed joint-resampling procedure showed material bias, so all of its confidence intervals and bootstrap uncertainty summaries are excluded from the active study. The excluded output is retained privately for audit and is not used as inferential support.
\section{Results}
\subsection{Primary associations}
Of 5,760 primary tests, 955 meet global $q<0.05$ (16.6\%): 557 have positive and 398 negative syndrome-minus-reference effects; 482 also have $|g|\geq0.8$. These are associations across correlated measurements, not independent anatomical discoveries. The complete supplementary table retains negative and nonsignificant findings.
\begin{figure}[htbp]\centering\includegraphics[width=\linewidth,height=.8\textheight,keepaspectratio]{figures/fig1_cross_syndrome_heatmap.png}\caption{Current syndrome--measurement associations. The heatmap displays twelve features selected by discovery frequency for description; inference retains the complete 5,760-test family.}\end{figure}
\begin{table}[htbp]\centering\small
\begin{tabular}{p{2.8cm}p{3.8cm}rp{2.7cm}r}\toprule Syndrome & FaceKit index & $g$ & 95\% CI & Global $q$\\\midrule
CLINICAL_ROWS
\bottomrule\end{tabular}
\caption{Prespecified examples. Intervals are marginal bootstrap intervals, while significance uses the complete primary BH family. Philtrum and ocular widths are normalized ratios; eyebrow area and malar geometry are indices rather than direct clinical diagnoses.}
\end{table}
\begin{figure}[htbp]\centering\includegraphics[width=\linewidth]{figures/fig2_clinical_indices.png}\caption{Syndrome-minus-reference Hedges' $g$ and marginal 95\% intervals for five prespecified FaceKit measurements.}\end{figure}
\subsection{Recorded-demographic associations}
\begin{table}[htbp]\centering
\begin{tabular}{llrr}\toprule Layer & Analysis & Hypotheses & BH associations\\\midrule
DEMOGRAPHIC_ROWS
\bottomrule\end{tabular}
\caption{Separate recorded-demographic correction families. Counts are not additive. Sparse support and source-label constraints limit interpretation.}\end{table}
\begin{figure}[htbp]\centering\includegraphics[width=\linewidth]{figures/fig3_reference_demographics.png}\caption{Reference-only omnibus associations for broad and detailed recorded groups, each with a separate 120-feature BH family.}\end{figure}
\subsection{Within-GMDB and adjusted analyses}
Against all other age-eligible GMDB patients, 222 associations pass BH; 173 overlap the 955 FairFace associations, with 782 FairFace-only and 49 within-GMDB-only findings. Restricting the comparator to other eligible syndromes gives 259 associations. The descriptive all-feature direction agreement against all other GMDB patients is 0.713 and Pearson effect correlation is 0.688; no concordance confidence intervals are reported.
The prespecified adjusted screen supports 5,640 of 5,880 planned models and identifies 507 BH associations. The targeted pairwise family supports 21,840 of 22,320 hypotheses and identifies 1,844 associations. The age-stratified family supports 2,760 of 11,760 hypotheses and identifies 39 associations. Unsupported rows and declared correction scopes are supplied alongside all supported estimates, including the separate four- and three-model covariate analyses.
\subsection{Reference and photograph robustness}
The all-eligible and category-by-recorded-sex reference specifications yield 971 and 945 BH associations, respectively. Deterministic alternate-photo selection changes 93 selected photographs and yields 996 associations: 918 shared, 37 primary-only and 78 alternate-only. One hundred random-photo replicates retain one image per patient and correct the complete family each time. Their per-hypothesis effect ranges and significance frequencies describe sampling sensitivity conditional on this cohort.
\begin{figure}[htbp]\centering\includegraphics[width=.85\linewidth]{figures/fig4_photo_selection.png}\caption{Primary and alternate-photo association status across all 5,760 hypotheses.}\end{figure}
\subsection{Separate synthetic catalog}
PDIDB revision \texttt{d0907a56525dc3b25712d360f9ce3031150eb535} contains 1,727 author-designated synthetic records; 1,563 pass QC. Model and verification tags do not establish generation lineage, patient-source independence or clinical validity. Age, sex and race fields are missing for all records. PDIDB supplies zero observations to primary, demographic, within-GMDB, adjusted or robustness inference. Image-publication permission remains unresolved; no synthetic image illustrations are used.
\section{Discussion and limitations}
PRISM describes statistical structure in FaceKit-derived geometric measurements and assesses how source, recorded groups, reference composition and photograph selection affect the associations. Within-GMDB comparisons address a different estimand from syndrome-versus-FairFace tests and do not establish clinical specificity. Adjustment and overlapping reference alternatives do not provide external replication.
Disease status remains confounded with GMDB-versus-FairFace source. Restoration, landmark errors, residual pose, expression and occlusion can affect geometry; normalized ratios can change through their denominators. Measurements are not calibrated three-dimensional anthropometry. Exact ages and clinical health are unavailable for FairFace, and unique reference identities are unknown. Sparse HC3 models and marginal bootstrap intervals may be unreliable. Concordance uncertainty is not reported because the assessed resampling method was biased. Correlated features and shared observations do not establish realized FDR. Independent, identity-aware samples with compatible acquisition and blinded geometric review are needed before diagnostic or clinical claims.
\section{Conclusion}
Using a single pinned set of 120 FaceKit measurements, PRISM identifies exploratory cohort associations with explicit source roles, correction families and robustness analyses. Its contribution is the statistical investigation; clinical interpretation and diagnostic utility require separate validation.
\section{Data, Code and Ethics}
Real GMDB photographs, patient identifiers, landmarks and per-person measurements remain restricted and are not supplied with manuscript materials. Dashboard distributions contain aggregate bins and summaries. Synthetic records remain separate from real-patient inference. The PI must confirm the exact IRB scope and identifier, image-publication permissions, authorship, funding and final release statement before submission. Provider review is needed for PDIDB generation lineage and redistribution terms. Existing repository-history exposure of earlier per-row data requires a separate PI remediation decision; preserving a private archive does not erase prior Git objects or remote copies.
\begin{thebibliography}{9}
\bibitem{facekit2026} H. Chen et al., ``FaceKit: a Toolkit for Interpretable Facial Phenotyping, Synthetic Image Generation and Privacy Analysis in Rare Diseases,'' arXiv:2610.09288v1, 2026.
\bibitem{gestaltmatcher} T.-C. Hsieh et al., ``GestaltMatcher facilitates rare disease matching using facial phenotype descriptors,'' Nature Genetics, vol. 54, pp. 349--357, 2022.
\bibitem{fairface} K. Karkkainen and J. Joo, ``FairFace: Face Attribute Dataset for Balanced Race, Gender, and Age for Bias Measurement and Mitigation,'' WACV, pp. 1548--1558, 2021.
\end{thebibliography}
\end{document}
